%------------------------------------------------------------------------------%
\section{Gráficos de Funções Trigonométricas}

Um fenômeno é chamado periódico se o mesmo se repete sempre após o mesmo
intervalo de tempo. Um exemplo simples de fenômeno periódico é o dia,
determinado pelo movimento do Sol. A observação do movimento periódico dos
corpos celestes nos possibilitou a criação do nosso calendário. Para uma melhor
compreensão dos fenômenos periódicos, faz-se necessário o estudo das funções
periódicas.

A repetição do valor numérico de uma função em um período determinado constitui
uma função periódica. Mais precisamente, uma função $f$ é periódica de período
$T$ se para todo $x$ real
\[
  f(x+T)=f(x).
\]

Observe que conhecido o gráfico de uma função periódica em um determinado
intervalo de seu domínio, que tenha o tamanho de seu período, automaticamente
conhecemos  o gráfico em todo o domínio, uma vez que podemos repetir o gráfico
em outros intervalos de mesmo tamanho. As funções trigonométricas são
periódicas e por isso apresentaremos os seus gráficos em intervalos do tamanho
de seus períodos. O período das funções seno e cosseno é $2\pi$. O período
da função tangente é $\pi$.

As seguintes observações merecem destaque.

As funções seno e cosseno são definidas em todo o conjunto dos
números reais e seus valores variam entre $-1$ e 1, ou seja, a imagem
dessas funções é o intervalo fechado de $[-1, 1]$. Isso significa que o
valor máximo que a função pode atingir é 1, enquanto o valor mínimo é
$-1$.

É importante notar que, além de possuir valor máximo de 1 e valor
mínimo de $-1$, as funções seno e cosseno atingem esses valores
infinitas vezes ao longo de seus domínios.

A função $y = \cos(x)$ é uma função par, o que significa que
\[
  \cos(-x) = \cos(x).
\]
Isso pode ser facilmente observado no círculo
trigonométrico, já que a função cosseno representa a coordenada $x$ de um
ponto na circunferência unitária, que é simétrica em relação ao eixo $y$.
Portanto, quando o ângulo é negativo, o cosseno mantém o mesmo valor
que quando o ângulo é positivo.

Para fazer o esboço do gráfico das funções $y=\sin(x)$ e
$y=\cos(x)$, utilizamos os arcos do eixo no círculo trigonométrico.

A função $y = \sin(x)$ é uma função ímpar, o que significa que
\[
  \sin(-x) = -\sin(x).
\]
Isso pode ser facilmente observado no círculo
trigonométrico.

Detalhes sobre função $y=\tan(x)$, será apresentado posteriormente,
no entanto, vale adiantar que essa função é ímpar, pois
\[
  \tan(-x)
  = \frac{\sin(-x)}{\cos(-x)}
  = \frac{-\sin(x)}{\cos(x)}
  = -\tan(x).
\]

As Figuras~\ref{fig:senono} e \ref{fig:cossenono} mostram o gráfico das funções
$y=\sin(x)$ e $y=\cos(x)$ em um período, ou seja, no intervalo $[0, 2\pi]$

%------------------------------------------------------------------------------%
\begin{figure}
  \begin{subfigure}[b]{0.5\linewidth}
    \centering
    \includegraphics[height=45mm]{seno-1}
    \caption{O gráfico da função seno em $[0, 2\pi]$.}
    \label{fig:senono}
  \end{subfigure}%
  \begin{subfigure}[b]{0.5\linewidth}
    \centering
    \includegraphics[height=45mm]{cos}
    \caption{O gráfico da função cosseno em $[0, 2\pi]$.}
    \label{fig:cossenono}
  \end{subfigure}
  \caption{Funções seno e cosseno.}
  \label{fig:sin_cos}
\end{figure}
%------------------------------------------------------------------------------%

Os próximos exemplos apresentam como construir gráficos de funções que são
transformações nas funções $y=\sin(x)$ e $y=\cos(x)$.

%------------------------------------------------------------------------------%
\begin{example}
  Esboce o gráfico da função
  \(f(x)= 2 \sin(x) +2\),
  explicitando o domínio e a imagem.

  \solution
  Note que $D(f)=\R$, uma vez que não temos nenhuma restrição para $x$.

  Para encontrarmos a imagem de $f$ devemos lembrar que
  \[
    -1 \leq \sin(x) \leq 1
  \]
  então se multiplicarmos por 2
  e depois somarmos 2 em todos os
  termos da desigualdade chegamos que
  \begin{align*}
    -2+2 & \leq 2 \sin(x) + 2 \leq 2+2 \\
    0    & \leq 2 \sin(x) + 2 \leq 4
  \end{align*}
  Assim, $\Im(f) = [0, 4]$.

  Uma vez que conhecemos o gráfico da função $y=\sin(x)$ podemos usar
  transformações para obtermos o gráfico de $f$. Uma vez que para obtermos
  a função $f$ multiplicamos a função $y=\sin(x)$ por 2 e depois somamos
  2 unidades, então para fazermos o esboço do gráfico de $f$ precisamos
  fazer uma expansão vertical de comprimento 2 e depois fazermos uma
  translação vertical (subimos 2 unidades no eixo $y$).

  \IncludeGraphicCentredEx{exemplo1seno}[0.5][-1][-1]

  Em azul o gráfico da função $y=2\sin(x)+2$
  e em vermelho gráfico da função $y=\sin(x)$.
\end{example}
%------------------------------------------------------------------------------%

%------------------------------------------------------------------------------%
\begin{example}
  Esboce o gráfico da função
  \(
   g(x) = \frac{1}{2} \cos\left( x + \frac{\pi}{2} \right) - 1
  \),
  explicitando o domínio e a imagem.

  \solution
  Note que $D(f)= \R$, uma vez que não temos nenhuma restrição para $x$.

  Para encontrarmos a imagem de $f$ devemos lembrar que
  \[
    -1 \leq \cos\left(x+ \frac{\pi}{2} \right) \leq 1
  \]
  então se multiplicarmos por
  \(\sfrac{1}{2}\) e depois subtraímos 1 em todos os termos dessa
  desigualdade chegamos que
  \begin{align*}
    -1-\frac{1}{2} & \leq \frac{1}{2} \cos\left(x+ \frac{\pi}{2}\right) -1 \leq \frac{1}{2} -1 \\
    - \frac{3}{2}  & \leq \frac{1}{2} \cos\left(x+ \frac{\pi}{2}\right) - \leq -\frac{1}{2}
  \end{align*}
  Assim, \(\Im(f)=\left[-\frac{3}{2},-\frac{1}{2}\right]\).

  Conhecemos o gráfico da função $y=\cos(x)$ então para obtermos o gráfico
  da função $f$ precisamos fazer três transformações na função $y=\cos(x)$:

  Contração vertical, por conta do fator $\sfrac{1}{2}$.

  Deslocamento horizontal, ou seja, ``andar para trás''
  \(\sfrac{\pi}{2}\) unidades.

  Fazer uma translação vertical, ou seja, ``descer'' uma unidade
  no eixo $y$.

  \IncludeGraphicCentredEx{exerciciocosseno}[0.5][-1][-1]

  Em azul, o gráfico da função
  \[
    y = \frac{1}{2} \cos\left( x + \frac{\pi}{2} \right) - 1
  \]
  e em vermelho, gráfico da função $y=\cos(x)$.
\end{example}
%------------------------------------------------------------------------------%

De modo geral, se temos uma função trigonométrica da forma
\[
  y = a + b \sin(kx+r)
  y = a + b \cos(kx+r)
\]
a imagem de qualquer uma dessas funções
é o intervalo $[a-b, a+b]$ e o período é
\[
  \frac{2\pi}{\vert k\vert}.
\]
No entanto, ao invés de olharmos só para essas fórmulas é interessante
entender o porquê disso.

A função $f(x) = \tan x$ não está definida entre
$-\sfrac{\pi}{2}$ e $\sfrac{\pi}{2}$,
e consequentemente (porque $f$ tem período $\pi$), em qualquer $x$
da forma
\[
  x = \frac{\pi}{2} + k\pi, \text{ para } k \in \Z,
\]
logo, o domínio da função tangente é
\[
  \left\{
    x \in \R \st x \neq \frac{\pi}{2+k\pi},
    \text{ para todo } k \in \Z
  \right\}
\]

\begin{figure}
  \centering
  \includegraphics[width=0.8\textwidth,trim={1cm 2cm 1cm 1cm},clip]{tg}
  \caption{O gráfico da função tangente em $(-\pi/2, \pi/2)$.}
\end{figure}

%------------------------------------------------------------------------------%